\documentclass[11pt]{article}

\usepackage[utf8]{inputenc}
\usepackage[T1]{fontenc}
\usepackage{amsmath}
\usepackage{amssymb}
\usepackage{array}
\usepackage{booktabs}
\usepackage{caption}
\usepackage{enumitem}
\usepackage{geometry}
\usepackage{graphicx}
\usepackage{microtype}
\usepackage{tabularx}
\usepackage{xcolor}
\usepackage{hyperref}

\geometry{margin=1in}
\hypersetup{
  colorlinks=true,
  linkcolor=blue,
  citecolor=blue,
  urlcolor=blue
}
\setlist{nosep}
\newcolumntype{Y}{>{\raggedright\arraybackslash}X}
\newcommand{\brfss}{\textsc{BRFSS}}
\newcommand{\gspan}{\textsc{gSpan}}
\newcommand{\phiboot}{P_{\mathrm{boot}}(\phi \ge 0.12)}

\title{Reproducible Multimorbidity Network Motifs Across Socioeconomic
Populations: Independent-Year Replication in BRFSS}
\author{}
\date{}

\begin{document}

\maketitle

\begin{abstract}
\textbf{Background:}
Multimorbidity is commonly summarized using disease prevalence, condition
counts, or pairwise associations, potentially obscuring recurrent
multi-disease structures across populations.

\textbf{Objective:}
To develop a survey-aware, bootstrap-stabilized framework for mining
multimorbidity network motifs and prospectively test their independent-year
reproducibility within the Behavioral Risk Factor Surveillance System
(\brfss).

\textbf{Methods:}
We used 457,670 \brfss{} 2024 records for discovery and 433,323 \brfss{} 2023
records for replication. Undirected graphs represented jurisdiction by
education-defined socioeconomic status (SES) strata, with 10 prespecified
chronic conditions as nodes. Disease-pair edges required a survey-weighted
point phi coefficient of at least 0.12 and, across 500 survey-design bootstrap
replicates, at least 90\% probability that phi remained at or above 0.12.
Connected, disease-labeled, non-induced motifs with three to five nodes were
mined from 94 frozen 2024 graphs. An independent bootstrap bank evaluated
motif robustness. Before accessing 2023 analytic data, we froze 484 robust
motifs as structural targets and three directional SES hypotheses as a
confirmatory family. No motif discovery was performed in 2023. Confirmatory
inference used density-residualized within-jurisdiction contrasts, 99,999
studentized Rademacher wild sign flips, and Holm correction.

\textbf{Results:}
The 2023 analysis included 47 paired jurisdictions, 94 graphs, and 951 stable
edges. Of 484 frozen motifs, 448 reproduced with at least 20\% pooled support
(92.6\%). Motif support was strongly correlated between years
($\rho=0.914$). All three focal motifs recurred structurally, but only one of
three SES hypotheses replicated. The lower-SES
arthritis--chronic-obstructive-pulmonary-disease--current-asthma--depression--
diabetes motif replicated (effect $=+0.2743$; one-sided $p=0.00105$; Holm
$p=0.00315$). Two higher-SES hypotheses were direction-concordant but did not
meet the frozen replication rule.

\textbf{Conclusions:}
A bootstrap-stabilized multimorbidity motif vocabulary showed strong
independent-year reproduction within \brfss, while SES-associated replication
was selective. These findings support external evaluation in distinct cohorts
and population-health data systems.
\end{abstract}

\noindent\textbf{Keywords:}
multimorbidity; network analysis; frequent subgraph mining; complex surveys;
bootstrap stability; socioeconomic status; replication; BRFSS

\section{Introduction}

Multimorbidity is often characterized using the prevalence of individual
conditions, the number of conditions per person, or pairwise disease
associations. These summaries are useful but may not capture higher-order
patterns in how disease associations are organized across populations.
Population-specific disease-association networks offer a complementary
representation: nodes denote chronic conditions, edges denote estimated
within-population associations, and recurring connected subgraphs describe
multi-disease configurations shared across populations.

Frequent subgraph mining can identify these recurrent configurations, but its
application to population health raises several methodological problems.
Hard association thresholds can admit sampling-sensitive edges; multi-edge
motifs may be less stable than their constituent edges; graph density can
mechanically inflate motif occurrence; and exploratory motif selection can
produce optimistic evidence if the same data are treated as confirmatory.
Complex survey data add the further requirement that weighting, strata, and
primary sampling units be respected when estimating associations and their
stability.

We addressed these problems in a discovery and replication design. First, we
developed and calibrated survey-weighted multimorbidity graphs using
\brfss{} 2024, stabilized their edges with a survey-design bootstrap, mined
recurrent motifs, and evaluated motif robustness and density sensitivity.
Second, before accessing 2023 analytic data, we froze the robust 2024 motif
vocabulary and three directional SES hypotheses. We then evaluated those
targets prospectively in \brfss{} 2023 without rerunning motif discovery.
\brfss{} 2023 therefore provides independent-year replication within the same
surveillance system, not external cohort validation.

\section{Methods}

\subsection{Data source and study populations}

\brfss{} is a repeated cross-sectional survey of noninstitutionalized adults
in the United States. The discovery analysis used 457,670 records from 2024,
and the replication analysis used 433,323 records from 2023. Survey design
variables included the final weight (\texttt{\_LLCPWT}), stratum
(\texttt{\_STSTR}), and primary sampling unit (\texttt{\_PSU}).

Each graph represented a jurisdiction by SES stratum rather than an
individual. Candidate jurisdictions were the 50 states and District of
Columbia; territories were excluded. A population stratum required at least
1,000 unweighted respondents, and both education strata had to be eligible for
a jurisdiction to enter paired analyses. The final discovery and replication
databases each included 47 paired jurisdictions and 94 graphs. Under the
frozen 2023 rules, District of Columbia and Nevada failed paired sample-size
eligibility, while Kentucky and Pennsylvania had no reporting data.

SES was defined using \texttt{\_EDUCAG}. The lower-education group comprised
respondents who did not graduate high school or graduated high school
(codes 1--2). The higher-education group comprised respondents who attended
college or technical school or graduated from college or technical school
(codes 3--4). Missing or unknown education values were excluded from the SES
domains. Population registries reported both unweighted sample size and Kish
effective sample size.

\subsection{Chronic-condition nodes}

The node set was frozen at 10 binary chronic-condition indicators:
ischemic heart disease (IHD), stroke, current asthma, skin cancer, other
cancer, chronic obstructive pulmonary disease (COPD), depressive disorder,
kidney disease, arthritis, and diabetes. Year-specific source variables and
coding rules were harmonized through a versioned condition registry. A node
required at least 500 valid observations and 30 cases within a population
stratum.

With 10 nodes, a complete undirected graph contains at most
\begin{equation}
  {10 \choose 2}=45
\end{equation}
disease pairs. The actual number of eligible dyads could be lower because each
pair required at least 500 complete observations, at least 30 cases of each
condition, and at least 10 unweighted co-occurrences.

\subsection{Survey-weighted disease-association graphs}

For binary disease indicators $A$ and $B$, we estimated a survey-weighted phi
coefficient within each jurisdiction by SES population. Let
$\widehat{p}_{11,w}$ denote weighted joint prevalence and
$\widehat{p}_{1+,w}$ and $\widehat{p}_{+1,w}$ denote weighted marginal
prevalences. Then
\begin{equation}
  \phi_w =
  \frac{
    \widehat{p}_{11,w}
    - \widehat{p}_{1+,w}\widehat{p}_{+1,w}
  }{
    \sqrt{
      \widehat{p}_{1+,w}(1-\widehat{p}_{1+,w})
      \widehat{p}_{+1,w}(1-\widehat{p}_{+1,w})
    }
  }.
\end{equation}
Phi is the Pearson correlation for two binary variables and is symmetric, so
the resulting graph was undirected. Geography and SES were graph metadata, not
nodes. An edge represented a sufficiently strong positive statistical
association in a population; it did not represent causation, direction,
disease progression, or a biological pathway.

\subsection{Survey bootstrap and frozen edge rule}

Edge stability was estimated using 500 survey-design bootstrap replicates.
The procedure did not independently resample respondents and did not resample
the 45 disease pairs. Instead, primary sampling units were resampled within
strata to form replicate survey weights. All eligible disease-pair phi
coefficients were recalculated under each replicate-weight realization, using
the same state-level replicate design for both SES domains.

For edge $e$, selection stability was
\begin{equation}
  \widehat{P}_{\mathrm{boot},e}(\phi \ge 0.12)
  =
  \frac{1}{500}
  \sum_{b=1}^{500}
  \mathbb{I}\!\left(\phi_e^{(b)} \ge 0.12\right).
\end{equation}
The final graph admitted an edge only when
\begin{equation}
  \boxed{
    \phi_{\mathrm{point}} \ge 0.12
    \quad\land\quad
    P_{\mathrm{boot}}(\phi \ge 0.12) \ge 0.90
  }.
  \label{eq:stable-edge}
\end{equation}
Every point-eligible dyad had to produce all 500 valid bootstrap estimates.
A failure made that population graph unusable and excluded the corresponding
paired jurisdiction.

\subsection{Calibration of the graph definition}

Graph calibration compared competing definitions rather than applying
sequential filters. Candidate rules included weighted phi thresholds of 0.08
and 0.12 and an age/sex-adjusted survey-logistic model of the form
\begin{equation}
  \operatorname{logit}\!\left\{P(B=1)\right\}
  =
  \beta_0+\beta_1 A+\boldsymbol{\beta}_2^\mathsf{T}\mathrm{age}
  +\beta_3\mathrm{sex}.
\end{equation}
The disease-pair odds ratio was $\exp(\beta_1)$. Positive adjusted
associations required a lower 95\% confidence bound above one and
within-graph Benjamini--Hochberg $q\le 0.05$. Forward and reverse disease
orientations were recorded because separate weighted logistic fits need not
be numerically identical. The adjusted rule served as a calibration and
concordance analysis; it was not part of the final edge filter.

The adjusted definition produced graphs that were too dense for useful
subgraph mining, with a median of approximately 31 of 45 possible edges.
The raw phi threshold of 0.12 produced substantially more moderate topology,
with approximately 18 edges per graph before stability filtering. These
results motivated Equation~\ref{eq:stable-edge} as the frozen graph
definition.

\subsection{Frequent subgraph mining}

The 94 frozen 2024 graphs were mined with \gspan. Motifs were connected,
disease-labeled, undirected, and non-induced, with three to five nodes.
Non-induced occurrence meant that all motif edges had to be present while
additional host-graph edges among motif nodes were allowed. Pooled support
thresholds were 10, 19, and 29 graphs, corresponding approximately to 10\%,
20\%, and 30\% of the 94-graph database. The primary discovery threshold was
20\%.

Support was opportunity-aware. A jurisdiction entered a motif's support
denominator only when every required motif dyad was eligible in both SES
graphs. At least 40 paired jurisdictions were required. Canonical disease and
edge order produced deterministic motif identifiers. To characterize
redundancy, we also identified occurrence-equivalence classes, closed motifs,
and maximal motifs.

\subsection{Motif robustness}

The primary robustness analysis used the frozen motif vocabulary and an
independent 500-replicate survey-bootstrap evaluation bank, distinct from the
bank that selected stable edges. Only edges already admitted to the frozen
stable graph were eligible in this analysis. A frozen edge could disappear
when its replicate phi fell below 0.12, but a previously rejected edge could
not enter. This estimand assessed joint stability of the graph structures
actually selected for analysis.

A secondary raw-threshold sensitivity allowed any eligible disease pair with
replicate phi at or above 0.12 to enter a replicate graph. This analysis
tested sensitivity to unrestricted threshold crossing and was not treated as
the primary motif-stability analysis.

\subsection{Density control and exploratory 2024 SES analysis}

Because denser graphs mechanically contain more subgraphs, motif occurrence
was standardized for graph-specific opportunity and edge count. For an
evaluable motif with $r$ required edges in a graph with $N$ eligible dyads and
$m$ selected edges, the fixed-edge-count expected occurrence probability was
\begin{equation}
  P_{\mathrm{null}}(\mathrm{occurrence})
  =
  \frac{{N-r \choose m-r}}{{N \choose m}},
  \qquad m\ge r.
\end{equation}
The density residual was observed occurrence minus expected occurrence. A
degree-preserving randomization constrained to eligible dyads served as a
sensitivity analysis.

Exploratory SES contrasts compared lower- and higher-education residuals
within jurisdiction. A common lower/higher label swap was applied within each
motif-evaluable state, and maxT adjustment controlled multiplicity across the
motif screen. These analyses described structural SES contrasts rather than
causal effects.

\subsection{Prospective 2023 replication}

The Phase 4 protocol, graph rules, 484-motif robust vocabulary, and three
directional SES hypotheses were frozen and tagged before 2023 analytic data
were accessed. The 2023 workflow reconstructed graphs under
Equation~\ref{eq:stable-edge}, evaluated every frozen motif without rerunning
\gspan, and retained all 484 motifs in the structural-reproduction
denominator. A motif reproduced when pooled support was at least 20\% among at
least 40 paired, motif-evaluable jurisdictions.

The confirmatory effect was the lower-minus-higher density-residual
difference. Thus negative effects favored the two higher-SES hypotheses and a
positive effect favored the lower-SES hypothesis. The test statistic was the
studentized, equally weighted mean state contrast. Primary inference used
99,999 common Rademacher sign vectors and one-sided alternatives in frozen
directions. Holm correction controlled the fixed three-hypothesis family at
$\alpha=0.05$. The sign-flip procedure is finite-sample exact only under
componentwise sign symmetry or within-state lower/higher exchangeability; its
stated justification was an asymptotic studentized wild-bootstrap argument
for a mean jurisdiction-level contrast. The matched 2024-jurisdiction subset
was descriptive and could not alter primary replication decisions.

\subsection{Reproducibility and audit trail}

The frozen protocol identifier was
\texttt{phase4\_2023\_cc5bbcc08283}. The original execution tag was
\texttt{phase4-2023-replication-v1.0.0}. The first 2023 execution stopped
during pre-harmonization runtime validation because YAML parsed Survey
package version 4.5 as a numeric scalar while the runtime returned the same
version as text. Deviation \texttt{PH4-D001} documented and authorized string
normalization for version comparison only. The original protocol remained
unchanged; no respondent values had been harmonized and no graph, motif, or
confirmatory result existed when the deviation was defined. The completed run
used execution tag \texttt{phase4-2023-replication-v1.0.1}. Tests, source
checks, manifests, and archive checksums were retained.

\section{Results}

\subsection{Edge stability and structural feasibility in 2024}

In 2024, 1,733 disease pairs met the raw point-phi threshold of 0.12. Of
these, 1,057 survived the 0.90 bootstrap-selection criterion, for 61.0\% raw
edge retention. The Phase 2.6 calibration protocol had prespecified an 80\%
retention criterion, so this result was formally recorded as
\textsc{HOLD} on retention. Instability was concentrated near the hard
threshold: 560 of 676 rejected raw edges had point phi below 0.15.

The retained graph database nevertheless remained structurally useful. Across
94 graphs, the mean stable edge count was 11.24 and the median was 11;
92 graphs had at least five edges, 90 had a connected component containing at
least five nodes, and 85 contained at least one triangle. The calibration
decision was therefore \textsc{HOLD} on the prespecified retention criterion
and \textsc{GO} on structural feasibility.

\subsection{2024 motif discovery and robustness}

\gspan{} identified 1,098 motifs with at least 10\% pooled support, 542 at the
primary 20\% support threshold, and 342 at 30\% support. The full discovered
set comprised 448 graph-occurrence-equivalence classes, 518 closed motifs,
and 114 maximal motifs. These counts describe a redundant computational
vocabulary rather than independent epidemiologic phenomena.

Among the 542 primary-support motifs, 499 had independent-bootstrap
$P_{\mathrm{boot}}(\mathrm{support}\ge20\%)\ge0.80$, and 484 (89.3\%) had
probability at least 0.90. Across frozen-edge \gspan{} reruns, median retention
of the baseline motif set was 95.0\%, with a 2.5th percentile of 88.7\%.
These 484 robust motifs, representing 299 robust occurrence families, became
the frozen 2023 structural-replication vocabulary.

The raw-threshold sensitivity produced a median of 4,999.5 motifs at 20\%
support, including a median of 4,457.5 motifs outside the frozen vocabulary.
Median motif-set Jaccard similarity with the frozen vocabulary was 0.108.
Thus unrestricted threshold crossing substantially altered the motif
landscape, supporting the stable-edge filter as a substantive part of the
analysis definition.

\subsection{Density diagnostics and exploratory 2024 SES signals}

Every retained motif had a positive support Z-score relative to the
eligible-dyad fixed-edge-count null, and 94.5\% had a positive Z-score
relative to the degree-preserving null. Because motifs had been selected from
the observed graphs, these null comparisons were selection-conditional
diagnostics rather than independent validation.

Exactly three motifs survived the exploratory paired maxT SES screen:
\begin{enumerate}
  \item \textbf{H1, higher SES:} an
  arthritis--IHD--diabetes--kidney--other-cancer tree, occurring in 12 lower-
  and 28 higher-SES graphs (maxT $p=0.0035$).
  \item \textbf{H2, higher SES:} an IHD--diabetes--kidney path, occurring in
  21 lower- and 35 higher-SES graphs (maxT $p=0.0167$).
  \item \textbf{H3, lower SES:} an
  arthritis--COPD--current-asthma--depression--diabetes tree, occurring in
  26 lower- and 11 higher-SES graphs (maxT $p=0.0321$).
\end{enumerate}
These were the only motifs from the exploratory screen that survived the
prespecified multiplicity procedure. Their structures and directions were
frozen prospectively for 2023.

\subsection{Independent-year structural replication in 2023}

The 2023 graph database included 47 paired jurisdictions, 94 graphs, and 951
stable edges. All point-eligible dyads completed the 500-replicate bootstrap.
Median selected edge count was 11 in lower-SES graphs and 9 in higher-SES
graphs; median eligible-dyad density was 25.0\% and 20.0\%, respectively.
Table~\ref{tab:year-summary} summarizes discovery and replication databases.

\begin{table}[htbp]
  \centering
  \caption{Discovery and independent-year replication summary.}
  \label{tab:year-summary}
  \begin{tabularx}{\textwidth}{Yrr}
    \toprule
    Measure & 2024 discovery & 2023 replication \\
    \midrule
    Raw annual records & 457,670 & 433,323 \\
    Paired jurisdictions & 47 & 47 \\
    Population graphs & 94 & 94 \\
    Stable edges & 1,057 & 951 \\
    Median stable edges per graph & 11 & 10 \\
    Motifs at primary 20\% support & 542 & not rediscovered \\
    Frozen robust motif targets & 484 & 484 evaluated \\
    Structurally reproduced motifs & not applicable & 448 (92.6\%) \\
    \bottomrule
  \end{tabularx}
\end{table}

Of the 484 frozen motifs, 448 reproduced with at least 20\% pooled support:
\begin{equation}
  \frac{448}{484}=92.6\%.
\end{equation}
The Spearman correlation between 2024 and 2023 pooled support was
\begin{equation}
  \rho=0.914.
\end{equation}
Median absolute support difference was 6.38 percentage points. All 42
three-node motifs reproduced, compared with 124 of 134 four-node motifs and
282 of 308 five-node motifs. Every motif had 46 or 47 evaluable paired
jurisdictions. All three focal confirmatory motifs met the structural
reproduction criterion in 2023.

\subsection{Confirmatory 2023 SES replication}

All three 2023 effects were concordant with their frozen directions, but only
H3 met the directional Holm-corrected replication rule
(Table~\ref{tab:confirmatory}).

\begin{table}[htbp]
  \centering
  \caption{Frozen SES hypotheses and prospective 2023 replication results.
  The 2024 lower/higher columns show raw graph occurrence; inference was
  density residualized.}
  \label{tab:confirmatory}
  \scriptsize
  \resizebox{\textwidth}{!}{
  \begin{tabular}{@{}lcrrrrrrl@{}}
    \toprule
    ID & Frozen direction & 2024 L/H & 2024 maxT $p$ &
    2023 support & 2023 effect & 2023 $p$ & Holm $p$ & Call \\
    \midrule
    H1 & Higher & 12/28 & 0.0035 & 34/94 & $-0.0449$ &
    0.26384 & 0.26384 & No \\
    H2 & Higher & 21/35 & 0.0167 & 47/94 & $-0.1705$ &
    0.05308 & 0.10616 & No \\
    H3 & Lower & 26/11 & 0.0321 & 27/94 & $+0.2743$ &
    0.00105 & 0.00315 & Yes \\
    \bottomrule
  \end{tabular}
  }
\end{table}

H1 remained direction-concordant but failed replication
(effect $=-0.0449$; one-sided $p=0.26384$; Holm $p=0.26384$). H2 also
remained direction-concordant but failed the frozen rule
(effect $=-0.1705$; one-sided $p=0.05308$; Holm $p=0.10616$). H2 was not
classified as ``near replication.'' H3 replicated
(effect $=+0.2743$; one-sided $p=0.00105$; Holm $p=0.00315$).

The descriptive 45-jurisdiction matched sensitivity analysis had the same
Holm-threshold pattern, with H3 alone meeting the statistical threshold, and
could not alter primary decisions. No aggregate study-level success gate was
preregistered. The confirmatory SES result was therefore one of three
hypotheses replicated.

\begin{figure}[htbp]
  \centering
  \fbox{\parbox{0.94\linewidth}{
    \textbf{Panel A: Analysis pipeline.}
    BRFSS respondents $\rightarrow$ jurisdiction by education-SES graphs
    $\rightarrow$ survey-weighted phi $\rightarrow$ survey bootstrap
    $\rightarrow$ stable edges $\rightarrow$ \gspan{} motifs
    $\rightarrow$ frozen 2024 targets $\rightarrow$ 2023 replication.
  }}

  \vspace{0.7em}
  \fbox{\parbox{0.94\linewidth}{
    \textbf{Panel B: Motif support replication.}
    Scatterplot of 2024 versus 2023 pooled support for all 484 frozen motifs,
    with the 20\% reproduction boundary, 448 reproduced motifs, and
    Spearman $\rho=0.914$.
  }}

  \vspace{0.7em}
  \fbox{\parbox{0.94\linewidth}{
    \textbf{Panel C: Confirmatory SES hypotheses.}
    Three motif diagrams with 2023 effect estimates and Holm-adjusted
    p-values, visually distinguishing the replicated H3 contrast from the
    non-replicated H1 and H2 contrasts.
  }}
  \caption{Proposed main figure for the discovery-to-replication study
  design and findings. Boxes are placeholders for final vector graphics.}
  \label{fig:main}
\end{figure}

\section{Discussion}

This study's principal methodological finding is that a survey-aware,
bootstrap-stabilized multimorbidity motif vocabulary reproduced strongly in a
different annual \brfss{} sample. More than 92\% of the 484 frozen motifs
again met the prespecified pooled-support criterion, and the support ordering
was highly consistent across years. This result indicates that recurrent
multimorbidity association structure was not confined to the 2024 discovery
sample.

The SES findings were narrower. All three focal structures recurred in 2023,
and their density-residualized effects retained the 2024 directions, but only
the lower-SES H3 contrast met the frozen Holm-corrected criterion. The H3
structure linked arthritis, COPD, current asthma, depressive disorder, and
diabetes. Its replication is compatible with a recurring concentration of
musculoskeletal, respiratory, mental-health, and metabolic association
structure in lower-education populations. However, the graph design cannot
determine whether this pattern reflects causal effects of education, shared
risk factors, health-care access, measurement, or other mechanisms.

The two higher-SES hypotheses illustrate the value of prospective
replication. Both remained direction-concordant, and H2 retained substantial
pooled structural support, but neither met its frozen confirmatory rule.
Treating H2 as a positive result because its unadjusted one-sided p-value was
close to 0.05 would violate the preregistered familywise decision procedure.
The appropriate conclusion is selective, one-of-three SES hypothesis
replication rather than universal confirmation of the exploratory 2024
screen.

The calibration analyses also show why representation stability matters.
A permissive or adjusted edge rule produced dense graphs, while unrestricted
replicate threshold crossing generated thousands of alternative motifs.
Requiring both a minimum point phi and 90\% bootstrap-selection stability
yielded a sparser graph database that remained structurally rich. The formal
Phase 2.6 retention failure remains important: many raw 0.12 edges were
sampling-sensitive. The graph database became useful not because the hard
threshold was intrinsically stable, but because the stability filter removed
many borderline edges.

\section{Limitations}

First, \brfss{} is self-reported and repeated cross-sectional, so disease
misclassification, nonresponse, and changes in survey administration may
affect estimated associations. Second, graph edges are statistical pairwise
associations, not causal links or evidence that an individual jointly has all
diseases in a motif. Third, the final phi graphs were not age/sex
standardized. The age/sex-adjusted logistic model was a competing calibration
definition and concordance check, not the final graph estimator.

Fourth, graph-specific density standardization addresses mechanical
subgraph-frequency differences but cannot eliminate all confounding or
differences in edge-detection precision caused by prevalence and sample size.
Fifth, edge-selection bootstrap stability does not fully propagate the
uncertainty of the entire graph-selection pipeline into confirmatory
p-values. Sixth, the Rademacher procedure relies on sign-symmetry for exact
finite-sample inference and otherwise on an asymptotic wild-bootstrap
argument.

Seventh, the 2023 analysis was independent-year replication within \brfss,
not validation in a different cohort or data-generating environment. Eighth,
the analysis included 47 paired jurisdictions per year and a prespecified
10-condition node set, which limits generalizability to other diseases and
geographies. Finally, frequent-subgraph outputs are nested and redundant;
raw motif counts should not be interpreted as counts of independent
epidemiologic phenomena.

\section{Conclusion}

Recurrent multimorbidity network structure was highly reproducible across
two \brfss{} years after survey-design bootstrap stabilization. One
lower-SES five-condition configuration showed confirmatory temporal
replication, while two higher-SES signals did not meet their frozen
familywise criteria. These findings support evaluating the frozen structures
in independent cohorts and other population-health data systems to test
portability beyond \brfss.

\section*{Data and code availability}

\brfss{} annual data are publicly available from the US Centers for Disease
Control and Prevention. Analysis code, frozen protocols, checksummed result
archives, and audit records are available at
\url{https://github.com/Efertugrul/multimorbidity-graphs}.

\appendix

\section{Supplementary calibration record}

The initial five-state prototype established feasibility but produced
relatively dense graphs. Expansion to all eligible states and District of
Columbia compared weighted phi thresholds and an age/sex-adjusted
survey-logistic graph definition. The adjusted definition had a median of
approximately 31 edges among 45 possible dyads and was retained only as a
calibration and concordance analysis. A raw phi threshold of 0.12 produced
approximately 18 edges per graph before stability filtering.

The 2024 raw-edge stability analysis retained 1,057 of 1,733 raw
phi-at-least-0.12 edges. Because this 61.0\% retention did not meet the
prespecified 80\% criterion, the calibration record was \textsc{HOLD}.
Structural diagnostics separately supported \textsc{GO}: nearly all graphs
had at least five edges and a component of at least five nodes, and 85 of 94
graphs contained triangles. This distinction was frozen before motif mining.

\section{Supplementary robustness and null-model interpretation}

The primary motif bootstrap conditioned on the frozen stable-edge universe:
selected edges could disappear, but rejected edges could not enter. This
estimand tested the joint robustness of motifs in the chosen graph
representation. The unrestricted raw-threshold sensitivity answered a
different question and showed a much less stable motif universe. Neither
analysis implemented a nested outer bootstrap that re-estimated the 0.90
selection-stability filter within every replicate.

Fixed-edge-count and degree-preserving null results were selection-conditional
because the motifs had been discovered in the observed graphs. They were used
to assess whether recurrence was trivially explained by graph density, not as
independent hypothesis tests. The strongest validation evidence was the
independent evaluation bootstrap followed by the prospectively frozen 2023
replication.

\section{Supplementary reproducibility identifiers}

\begin{tabularx}{\textwidth}{@{}lY@{}}
  \toprule
  Item & Identifier \\
  \midrule
  Frozen Phase 4 protocol & \texttt{phase4\_2023\_cc5bbcc08283} \\
  Frozen protocol release & \texttt{phase4-2023-replication-v1.0.0} \\
  Documented execution deviation & \texttt{PH4-D001} \\
  Completed execution release & \texttt{phase4-2023-replication-v1.0.1} \\
  2023 configuration digest & \texttt{1f0234d0e0ed} \\
  Result archive & \texttt{results/baselines/phase4\_1f0234d0e0ed} \\
  Result release & \texttt{phase4-2023-results-v1.0.0} \\
  \bottomrule
\end{tabularx}

\end{document}
